%==============================================================================%
% Chapter 10 – Jasmin-to-EasyCrypt Refinement Proof
%==============================================================================%

\chapter{Jasmin-to-EasyCrypt Refinement Proof}
\label{ch:refinement}

\section{Overview of \texttt{RefJasminNTT.ec}}

The file \texttt{RefJasminNTT.ec} (1000+ lines) bridges the gap between the
machine-level \jasmin{} implementation and the abstract algorithmic
specification.  It is the most technically demanding proof file in the
development, requiring simultaneous reasoning about:
\begin{enumerate}
  \item Signed vs.\ unsigned word arithmetic.
  \item Array representation correspondences.
  \item Integer arithmetic bounds.
  \item Equivalence of three-level nested loops.
\end{enumerate}

\section{Word-to-Coefficient Correspondence}

\subsection{The Linking Predicate}

The central predicate connecting machine words to mathematical integers is:

\begin{lstlisting}[language=EasyCrypt,
  caption={Array correspondence predicate in \texttt{RefJasminNTT.ec}},
  label={lst:wrel}]
pred wrel (rp : W32.t Array256.t) (a : int Array256.t) =
  forall i, 0 <= i < 256 =>
    W32.to_sint rp.[i] = a.[i].
\end{lstlisting}

This predicate asserts that the 32-bit word representation of each element
in \texttt{rp} equals the integer in \texttt{a}.  Note that \ec{W32.to\_sint}
applies the two's-complement signed interpretation.

\subsection{Preservation Under Arithmetic}

\begin{lemma}[Word addition preservation]
  If $\ec{wrel}(\texttt{rp}, \texttt{a})$ and $|a[i]|, |b[i]| < 2^{30}$,
  then $\ec{wrel}(\texttt{rp} \oplus_{32} \texttt{b}, \texttt{a} + \texttt{b})$
  where $\oplus_{32}$ denotes 32-bit word addition.
\label{lem:wrel_add}
\end{lemma}

\begin{proof}
  The key identity is: for $a, b \in \Z$ with $|a|, |b| < 2^{30}$,
  \[
    \op{W32.to\_sint}(\op{W32.of\_int}(a) + \op{W32.of\_int}(b)) = a + b,
  \]
  which holds because the 32-bit arithmetic does not overflow when both
  operands are bounded by $2^{30}$.
\end{proof}

\section{Montgomery Reduction Correspondence}

\subsection{Word-Level Reduction}

The \jasmin{} procedure computes Montgomery reduction on machine words.
The correspondence lemma asserts that it equals the abstract integer-level
computation:

\begin{lemma}[Montgomery reduction word correspondence]
  Let $z \in \Z$ with $|z| < 2^{31} \cdot q$.  Then:
  \[
    \op{W32.to\_sint}(\op{Extracted.montgomery\_reduce}(\op{W64.of\_int}(z)))
    = \op{NTT\_Fq.montgomery\_reduce}(z).
  \]
\label{lem:mont_word_corr}
\end{lemma}

\begin{proof}
  The extracted procedure performs the same arithmetic as the integer-level
  specification.  The key step is showing that intermediate values do not
  overflow: the 64-bit product $m \cdot q$ satisfies $|mq| \le (2^{32}-1)
  \cdot 64513 < 2^{63}$, fitting in a signed 64-bit register.
\end{proof}

\section{Butterfly Equivalence}

\subsection{Individual Butterfly}

\begin{lemma}[Single butterfly correspondence]
  Let $\texttt{rp}, \texttt{a}$ satisfy $\ec{wrel}(\texttt{rp}, \texttt{a})$,
  and let $|a[j]|, |a[j+\ell]| < 2^{15}$ and $|\zeta| < 2^{15}$.  Then after
  executing one butterfly step:
  \[
    \ec{wrel}(\texttt{rp}', \texttt{a}'),
  \]
  where $\texttt{a}'$ is the result of the abstract butterfly and
  $\texttt{rp}'$ is the result of the machine butterfly.
\label{lem:butterfly_corr}
\end{lemma}

\begin{proof}
  The machine butterfly computes:
  \begin{align*}
    t &= \op{W32}(\op{Extracted.mont\_reduce}(\zeta \cdot \texttt{rp}[j+\ell])) \\
    \texttt{rp}'[j+\ell] &= \texttt{rp}[j] - t \\
    \texttt{rp}'[j] &= \texttt{rp}[j] + t
  \end{align*}
  By \Cref{lem:mont_word_corr}, $\op{W32.to\_sint}(t) = \op{mont}(\zeta \cdot a[j+\ell])$.
  By \Cref{lem:wrel_add}, the addition/subtraction preserves the correspondence
  (using bounds $|a[j]|, |\op{mont}(\zeta \cdot a[j+\ell])| < 2^{16}$, so
  their sum/difference is bounded by $2^{17} < 2^{30}$, safe in 32-bit words).
\end{proof}

\section{Loop Equivalence Proofs}

\subsection{Inner Loop Equivalence}

\begin{lemma}[Inner loop correspondence]
  Given $\ec{wrel}(\texttt{rp}, \texttt{a})$ and bounds $|a[i]| < 2^{15}$:
  \[
    \ec{hoare}[\ec{Extracted.inner\_loop} \sim \ec{NTT\_Fq.inner\_loop} :
      \ec{wrel}(\texttt{rp}, \texttt{a}) \Rightarrow \ec{wrel}(\texttt{rp}', \texttt{a}')].
  \]
\label{lem:inner_corr}
\end{lemma}

\begin{proof}
  By induction on the inner loop counter $j$, using \Cref{lem:butterfly_corr}
  as the inductive step.  The bound on coefficients increases by at most
  $2^{16}$ per butterfly (from \Cref{thm:mont_correct}), but the inner loop
  uses the same twiddle factor throughout, so the number of butterflies per
  inner loop execution is exactly $\ell$.
\end{proof}

\subsection{Full NTT Loop Equivalence}

\begin{theorem}[Forward NTT loop equivalence]
  Let $\texttt{rp}$ satisfy $|a[i]| < 2^{15}$ for all $i$.  Then:
  \begin{align*}
    &\ec{hoare}[\ec{Extracted.poly\_ntt\_jazz} \sim \ec{NTT\_Fq.NTT.ntt} : \\
    &\quad \ec{wrel}(\texttt{rp}, \texttt{a}) \\
    &\quad \Rightarrow \ec{wrel}(\texttt{rp}', \texttt{a}') \;\wedge\;
     \forall i.\; |a'[i]| < 2^{23}].
  \end{align*}
\label{thm:ntt_loop_equiv}
\end{theorem}

\begin{proof}
  By three-level loop induction:
  \begin{enumerate}
    \item The outer loop induction uses $\ell$ as the decreasing variant.
          The invariant after completing all stages with $\ell' > \ell$ is
          the partial NTT condition.
    \item The middle loop induction uses $256 - \texttt{start}$ as the
          decreasing variant.
    \item The inner loop uses \Cref{lem:inner_corr}.
  \end{enumerate}
  The bound $|a'[i]| < 2^{23}$ follows from \Cref{thm:ntt_bounds}.
\end{proof}

\section{Main Refinement Lemmas}

The two main lemmas of \texttt{RefJasminNTT.ec} are:

\begin{lstlisting}[language=EasyCrypt,
  caption={Main refinement lemmas in \texttt{RefJasminNTT.ec}},
  label={lst:ref_lemmas}]
(* Forward NTT: Jasmin core ~ NTT_Fq, bounds 16 -> 24 *)
lemma poly_ntt_core_ref rp a :
  wrel rp a =>
  (forall i, 0 <= i < 256 => `|a.[i]| < 2^15) =>
  equiv[Extracted.poly_ntt_jazz ~ NTT_Fq.NTT.ntt :
    ={arg} /\ wrel arg{1} arg{2} ==>
    wrel res{1} res{2} /\
    forall i, 0 <= i < 256 => `|res{2}.[i]| < 2^23].

(* Inverse NTT: Jasmin core ~ NTT_Fq, bounds 16 -> 16, Montgomery *)
lemma poly_invntt_core_ref rp a :
  wrel rp a =>
  (forall i, 0 <= i < 256 => `|a.[i]| < 2^23) =>
  equiv[Extracted.poly_invntt_jazz ~ NTT_Fq.NTT.invntt :
    ={arg} /\ wrel arg{1} arg{2} ==>
    wrel res{1} res{2} /\
    forall i, 0 <= i < 256 => `|res{2}.[i]| < 2^15].
\end{lstlisting}

These two lemmas represent the central contribution of \texttt{RefJasminNTT.ec}.
Once proved, they serve as the foundation for the end-to-end theorem.

\section{Proof Engineering Challenges}

\subsection{Signed vs.\ Unsigned Arithmetic}

The \jasmin{} code uses \jasm{u32} (unsigned) for coefficient storage but
\jasm{(64s)(x)} (signed 64-bit) for intermediate products.  The
\easycrypt{} extraction mirrors this by using \ec{W32.t} for storage and
\ec{W64.t} for products.  The refinement proof must track the signed
interpretation of unsigned words, which requires the signed-arithmetic
lemmas from \texttt{Fq.ec}.

\subsection{64-bit Product Non-Overflow}

The 64-bit product $\zeta \cdot a[j+\ell]$ requires:
$|\zeta| \cdot |a[j+\ell]| < 2^{63}$.
Since $|\zeta| < q < 2^{16}$ and $|a[j+\ell]| < 2^{23}$ (after the forward
NTT), we have $|\zeta \cdot a[j+\ell]| < 2^{39} \ll 2^{63}$.
This is verified in the proof as a side condition.

\subsection{Twiddle Factor Table Agreement}

The \easycrypt{} extracted module accesses the twiddle table as a global
constant.  The refinement proof must show that this global constant agrees
with the \ec{zetas} array defined in \texttt{NTT\_Fq.ec}.  This is
established by a computation lemma:

\begin{lemma}[Twiddle table agreement]
  For all $k \in [0, 256)$:
  $\op{W32.to\_sint}(\ec{jzetas}[k]) = \ec{NTT\_Fq.zetas}[k]$.
\label{lem:zetas_agree}
\end{lemma}

\begin{proof}
  Both tables are defined by the same formula modulo $q = 64513$.  Since
  $q < 2^{16}$, all entries fit in 16 bits and the signed interpretation
  agrees with the unsigned value.  Verified by the SMT backend.
\end{proof}

\subsection{In-Place Update Semantics}

Both the \jasmin{} implementation and the \easycrypt{} specification modify
the input array in place.  The refinement proof tracks the state of each
array cell individually, using the \easycrypt{} array theory lemmas to reason
about functional updates.

The key identity is: after an in-place butterfly that modifies indices $j$ and
$j+\ell$, all other indices remain unchanged.  This is expressed as:
\[
  \forall i \ne j, i \ne j+\ell.\;
  \texttt{rp}'[i] = \texttt{rp}[i] \;\wedge\;
  \texttt{a}'[i] = \texttt{a}[i].
\]
